\section{Conclusion}

Après délinéarisation, l'application de ces étapes sur la fonction \texttt{sum}
produit alors le code suivant, qui est bien récursif terminal :

\small
\begin{minted}{ocaml}
let rec sum_acc (l : int list) (acc : int) : int =
    match l with
    | [] -> acc
    | x :: q ->
        let new_acc (a_7 : int) : int = x + a_7 in
        sum_acc q (new_acc acc)

let sum (l : int list) : int = sum_acc l 0
\end{minted}

\normalsize

Les étapes décrites ci-dessus peuvent être appliquées à des fonctions plus
compliquées, par exemple des couples de fonctions mutuellement récursives ou des
fonctions effectuant plusieurs appels récursifs.

Finalement, l'application de la transformation implémentée par mon camarade
transforme ces fonctions récursives terminales en des fonctions implémentées
avec des boucles, donnant un résultat similaire à celui montré en exemple dans
la section 1.
